Beyond nitrogen and phosphorus, subsurface drainage governs the field-to-stream
transport of pesticides, manure-borne microbes and resistance genes, trace
elements, and dissolved greenhouse gases, and it reshapes agricultural ecology.
Across these dimensions the evidence converges on a common logic: the same
preferential-flow pathways and hydrologic controls that dominate nutrient loss
also govern contaminant delivery, and management---application method, timing,
and edge-of-field engineering---frequently matters more than drainage per se.

\subsection{Herbicides, Pesticides, and Emerging Contaminants}

The timing of runoff and drainage relative to application is the single most
powerful control on peak pesticide concentrations and losses. Synthesising more
than thirty North American studies, \textcite{Kladivko2001} established that
pesticides are commonly detected in drainflow during the first drainage event(s)
after application, concentrations decline rapidly thereafter, and most mass
losses are very small---generally less than $0.5\%$ and often below $0.1\%$ of
applied---while tile concentrations and losses are usually much lower than in
surface runoff, frequently by an order of magnitude. Site-specific values
corroborate this: Illinois long-term atrazine losses were $0.11$--$0.22\%$ of
application \parencite{Kalita2006}; New York 2,4-D and atrazine losses were
$0.08\%$ and $0.02\%$ \parencite{Shalit1996}; and less than $1\%$ of applied
herbicide was recovered in tile drains over two years
\parencite{Fortin2002,Siimes2006}. Early events dominate: most atrazine and
metolachlor loss occurred within the first $61$ days \parencite{SOUTHWICK1990},
and over $80\%$ of catchment losses occurred during just two rain events
\parencite{Leu2004}. In cold regions most loss can occur during snowmelt and
thaw rather than immediately after application \parencite{Siimes2006}, and
two-phase seasonal behaviour is common \parencite{Novak2001,Dousset2004}.

The mechanism reconciling near-immediate arrival with small overall losses is
preferential (macropore) flow. Tracer breakthrough curves were typical of
preferential flow, with herbicides arriving after less than $1$~cm of net
drainage, interpreted as transport through only $\sim3\%$ of pore volume
\parencite{Fortin2002,Dousset2004}. Soil structure, rather than sorption or
degradation, dominated between-soil differences, with greater metolachlor losses
in a clayey Vertic Cambisol than a loamy Luvisol \parencite{Novak2001,Novak2003}.
Although mass losses are small, transient concentrations can exceed thresholds:
Ontario maxima reached $287$~$\mu$g~L$^{-1}$ (atrazine) in tile discharge, with
the atrazine MCL exceeded in $60\%$ of tile samples \parencite{GAYNOR1995}, and
Iowa subsurface drains had glyphosate maxima above $250$~$\mu$g~L$^{-1}$
\parencite{Coupe2012}. Degradates contribute meaningfully and can be more mobile
than parents: deethylatrazine accounted for $10$--$27\%$ of combined loss
\parencite{GAYNOR1995,MUIR1978,JAYACHANDRAN1994}. Management levers include
banding, residue-conserving tillage with intercrops (reducing loss
$46$--$49\%$), and application-timing restriction via a soil-wetness index
\parencite{Gaynor2001,Julien2023,Moorman1999}, whereas corn-stover harvest can
worsen drainage quality \parencite{Shipitalo2016}. Woodchip and biochar-amended
bioreactors removed atrazine (up to $93\%$) mainly abiotically without inhibiting
denitrification \parencite{Hassanpour2019,Ilhan2011}, and modelling frameworks
(PESTDRAIN, HYDRUS-2D, RZWQM, SWAT+) reproduce drainage pesticide dynamics when
macroporosity and sorption kinetics are represented
\parencite{Branger2009,Boivin2006,Wendell2024,Novak2003}. Tile drainage also
transports emerging contaminants---PPCPs, hormones, PBDEs, perfluoroalkyl acids,
and neonicotinoids---whose behaviour mirrors the timing and preferential-flow
patterns above, generally at low, transient concentrations
\parencite{Gottschall2012,Gottschall2013,Gottschall2017,Ghirardini2019,Schaafsma2019}.

\subsection{Manure-Borne Microbial Contaminants and Resistance Genes}

Manure and organic amendments introduce fecal microorganisms, resistance genes,
hormones, and trace elements that drainage can rapidly route to receiving waters,
with application method and soil moisture---rather than manure presence
alone---largely governing transport. Curli-producing \textit{E.\ coli} were far
more abundant in tile drains from farms that surface-applied manure than one that
incorporated it, indicating that application method itself shapes the transported
phenotype \parencite{Truhlar2015}. Tillage governs pathogen delivery mainly
through macroporosity: zero tillage generally produced higher subsurface flow and
greater \textit{E.\ coli} loads than conventional tillage
\parencite{Thiagarajan2007,Jenkins2009}. Antecedent soil moisture and water-table
position exert first-order control: effluent on wet soil produced drain-flow
\textit{E.\ coli} of $10^{5}$--$10^{7}$ per~100~ml versus $\sim10^{4}$ under a
greater moisture deficit, and where preferential flow occurred, drainage
concentrations approached those of the applied effluent
\parencite{Collins2007,Ding2021}. Manure is also a reservoir of antibiotic
resistance genes that persist through drained soils into drainage water: shotgun
metagenomics identified $320$ unique ARGs, with manure treatment shifting soil
and drainage-water resistomes and many manure-derived genes associated with
mobile genetic elements in food- and water-borne pathogens \parencite{Smith2019b}.
Steroid hormones were detected only transiently and generally did not differ
significantly among treatments \parencite{Jenkins2009}. Engineered receiving
features can attenuate microbial loads---a constructed wetland reduced nitrate and
\textit{E.\ coli} mass loads $67.6\%$ and $63.3\%$ \parencite{Haverstock2017}---
though some wetlands act as net \textit{E.\ coli} exporters \parencite{Stott2023}.
Drained soils, bioreactors, and wetlands host distinct, dynamic microbial
communities structured by depth and season, and favourable conditions can be
engineered to develop effective denitrifier assemblages
\parencite{Andrus2014,Porter2015,Hathaway2015,Knorr2023,Schilling2017}.

\subsection{Trace Elements, Selenium, and Acid Sulphate Hazards}

Where drainage concentrates trace elements, it can create acute ecological
hazards. The Kesterson Reservoir remains the definitive case: waterborne selenium
averaged $0.122$~mg~L$^{-1}$ ($\sim230$ times reference), biomagnified through the
food chain, and caused dead and deformed bird embryos, with integrated studies
fulfilling the equivalent of Koch's postulates and identifying selenium as the
sole causative agent \parencite{SAIKI1987,Ohlendorf2002,BENSON1993}. The
regulatory legacy is durable: open-water evaporation is now prohibited in
California, requiring rock-media systems \parencite{Ayars2014,BRADFORD1991}, and
selenium-removal wetlands retained about $59\%$ of inflow Se \parencite{Gao2003}.
In evaporation basins, evapoconcentration accumulates extreme salinity and trace
elements \parencite{WESTCOT1988}, and redox chemistry jointly controls arsenic
and uranium mobility \parencite{Ryu2010,Gao2007,Duff1997}; boron isotopes and
other tracers can distinguish contaminant sources
\parencite{Tayfur2010,Komor1997,Zhou2023}. At the opposite extreme, intensive
drainage can drive negative micronutrient balances and net soil depletion
\parencite{Gul2013}, while central-Iowa drainage ionic composition was dominated
by bicarbonate, calcium, and magnesium \parencite{Zimmerman2017b}. Drainage of
acid sulphate soils is a distinct hazard: oxidation of pyritic sediments drives
stream pH declines to near or below $4.0$ and elevates aluminium and metals, with
controlled pipe drainage recommended over in-stream liming
\parencite{Osterholm2005,Åström2005,Virtanen2017,Salo2023}.

\subsection{Greenhouse-Gas Dimensions}

Drainage nitrate is coupled to greenhouse-gas fluxes, and drainage water is
commonly supersaturated with dissolved N$_2$O, CH$_4$, and CO$_2$. In Japan,
indirect emissions via subsurface drainage were equivalent to $4.6\%$, $58.3\%$,
and $2.5\%$ of direct emissions, with N$_2$O-N/NO$_3$-N ratios below the IPCC
factor \parencite{Sawamoto2003,Minamikawa2010,Minamikawa2013}. In Canada,
dissolved N$_2$O losses approached surface emissions, were strongly seasonal, and
declined along tile length, implying an in-drain sink and likely underestimation
of the IPCC coefficient \parencite{Burton2024}; a global synthesis found
subsurface drains had the highest mean emission factor among artificial waters,
with the stationary IPCC factor tending to overestimate N$_2$O production
\parencite{Webb2021}. Conversely, installing drainage sharply reduces methane in
paddy systems by improving aeration: drained fields emitted $71\%$ less CH$_4$
than non-drained fields \parencite{Shiratori2007,Furukawa2008}. Management
modulates N$_2$O: controlled drainage kept cumulative N$_2$O below $1\%$ of
applied N, and saturated riparian buffers removed nitrate without increasing
N$_2$O \parencite{Smith2019,Davis2019,Nash2015e,Carstensen2021}. Carbon export
from drained fields is dominated by inorganic carbon through tiles ($\sim96\%$ of
total carbon export), so reducing water throughput supports carbon retention
\parencite{Schilling2023,Driscoll2025}.

\subsection{Ecological Responses to Subsurface Drainage}

Drainage reshapes agricultural ecology from soil fauna upward. \textit{Lumbricus
terrestris} abundance was typically higher above drains than between them (median
counts twice as high, fresh mass five times), so drainage pattern affects
population density and should inform sampling design
\parencite{Nuutinen2001,Nuutinen2011}. Drainage type is associated with
farmland-bird outcomes, with skylark and meadow pipit favouring open-drained over
subsurface-drained fields \parencite{Marja2013}, while tile-drainage wetlands
carried cumulative pesticide concentrations $4.7$ times higher than reference
wetlands (neonicotinoid-dominated), with measurable neurological effects in
juvenile amphibians \parencite{Campbell2023}. Subsurface drainage of meadows was
a leading cause of decline in semi-aquatic snails \parencite{Dreyfuss2018}.
Restored aquatic habitats can benefit: in restored Iowa oxbows, tile drainage had
negligible effects on physical or biotic variables, providing promising evidence
that multipurpose oxbows can enhance biodiversity without ecological penalty, with
outcomes improved by sediment excavation
\parencite{Leberg2024,Winikoff2020,Harris2018}. Crucially, management practices
often exert clearer effects than tile drainage itself: no-till and diverse
rotations altered bulk density, organic matter, and aggregate stability while
cover crops and drainage had minimal effect in some settings
\parencite{Saha2024,Hermawan1996,Sherbine2023}, and additional practices such as
polyacrylamide amendment and anti-weed mesh with drainage provide further
quantified nutrient-retention and production benefits
\parencite{Jiang2010,GiménezValero2024,GiménezValero2025}.